\section{Related Work}

sEMG has long been studied as a non-invasive control channel for rehabilitation, prosthetics, and human-machine interaction, and stroke-specific reviews show that sEMG-driven interventions can support meaningful upper-limb rehabilitation when signal interpretation is reliable enough for assistive use~\citep{munoznovoa2022upper,zheng2021semg}. Broader reviews of wearable sensing for stroke recovery make a related point: clinical value depends on the full stack, including acquisition quality, body placement, labeling protocol, feature representation, and the way machine learning is coupled to assessment or assistive control~\citep{boukhennoufa2022wearable}. For this reason, decoding architecture should not be discussed independently of data construction and deployment context.

Dataset choice is especially important in this area. Ninapro remains the dominant public benchmark family for hand-gesture decoding from sEMG, with ten sub-datasets spanning different electrode layouts, movement vocabularies, and sensor modalities~\citep{ninapro2014,ninaproSite2026}. The widely used DB5 split, for example, records 52 movements from 10 healthy participants with two Myo armbands, while DB1 and DB2 cover larger healthy cohorts with different channel configurations and richer kinematic metadata~\citep{ninapro2014,ninaproSite2026}. By contrast, \physiomio\ was introduced in 2026 as a stroke-specific bilateral and longitudinal HD-sEMG resource with 48 stroke patients, 16 gestures, 64 electrodes, and repeated recordings across inpatient recovery~\citep{physiomioData2026}. That distinction matters because results on healthy-subject multiclass gesture sets are not automatically transferable to impaired-arm intent decoding in post-stroke populations.

Stroke-specific decoding studies have historically been small and heterogeneous. Lee et al.~\citep{lee2010stroke} used subject-specific myoelectric pattern classification for six functional hand movements in chronic stroke survivors and reported a sharp impairment-dependent gap, with mean accuracy of 71.3\% for moderately impaired participants and 37.9\% for severely impaired participants. Anastasiev et al.~\citep{anastasiev2022stroke} later studied post-acute stroke gesture recognition with portable eight-channel sEMG and found that affected-side SVM performance could still reach the high-80\% range on reduced gesture sets, but the task remained sensitive to feature design, label count, and paretic signal quality. These studies underscore a recurring challenge for post-stroke intent decoding: strong performance is possible, but dataset scale and protocol differences make general claims difficult.

Deep neural models broaden that picture but do not remove the comparability problem. Sequence-oriented recurrent models remain natural baselines because they aggregate temporal evidence across windows, while compact 1D CNNs emphasize local temporal motifs and can be compressed aggressively for efficient inference~\citep{hochreiter1997lstm}. Graph neural networks offer a complementary structural view by representing windows as related observations connected through message passing rather than as a purely linear sequence~\citep{kipf2017gcn,hamilton2017graphsage,velickovic2018gat}. On healthy-dataset benchmarks such as Ninapro, deeper end-to-end networks can reach very high multiclass accuracies; for example, Sri-Iesaranusorn et al.~\citep{sri2021dnn} reported 93.87\% accuracy on Ninapro DB5 and 91.69\% on DB7 for 41-movement classification. Those results are useful context for model capacity, but they address a different population and a different task than impaired-arm multilabel finger-intent prediction.

Transfer learning has recently become one of the most relevant directions for rehabilitation-oriented sEMG. Li et al.~\citep{li2025transfer} showed that transfer-based CNNs can materially improve inter-subject and inter-day robustness on high-density healthy-subject sEMG, and Mohammadiazni et al.~\citep{mohammadiazni2026transfer} demonstrated an especially important stroke-specific result: pretraining on healthy Ninapro DB5 data raised stroke-survivor gesture-classification accuracy from 46\% to 93.6\% in a three-gesture setting. These findings motivate the healthy-to-impaired transfer pathway in \repo, but they also raise a cautionary point. Transfer-learning gains are task-, architecture-, and dataset-dependent, so a project paper should report them as empirical outcomes rather than as guaranteed improvements.

Deployment-oriented optimization provides the final piece of context for this manuscript. Knowledge distillation is a standard way to transfer behavior from larger teachers into smaller students~\citep{hinton2015distilling}, and Optuna is widely used for black-box hyperparameter search under practical constraints~\citep{akiba2019optuna}. What distinguishes \repo\ is therefore not a novel recurrent, convolutional, or graph formulation in isolation, but an integrated software research effort that places dataset harmonization, baseline comparison, teacher-student modeling, transfer learning, and deployment proxies inside one auditable experimental framework.

\begin{table}[t]
  \centering
  \footnotesize
  \caption{Literature context used to position the current study. Reported numbers come from different populations, sensors, gesture vocabularies, and label spaces, so the table is for context rather than for a direct leaderboard comparison.}
  \label{tab:literature-context}
  \resizebox{\linewidth}{!}{\begin{tabularx}{\linewidth}{>{\RaggedRight\arraybackslash}p{2.15cm}YYYY}
\toprule
Source & Population / dataset & Task & Reported performance & Relevance to this paper \\
\midrule
Lee et al. 2010 & 20 chronic stroke survivors; 10 forearm/hand electrodes & Six functional hand movements with subject-specific EMG classification & Mean accuracy 71.3\% for moderate impairment and 37.9\% for severe impairment & Early evidence that stroke severity strongly affects decoder quality \\
Anastasiev et al. 2022 & 19 post-acute stroke patients; portable eight-channel sEMG & Four- to seven-gesture multiclass recognition on affected and non-affected sides & Affected-side SVM accuracy reached 88.7\% on the four-gesture setting & Shows that reduced-vocabulary stroke gesture recognition can reach high accuracy with careful feature design \\
Bao et al. 2024 & 8 chronic stroke participants; post-stroke sEMG feature-domain study & CNN, CNN-LSTM, and CNN-LSTM-attention recognition using time, frequency, and wavelet features & Best reported averages: 72.95\% intra-subject accuracy and 68.38\% inter-subject transfer accuracy & Supports deep architectures for post-stroke decoding while highlighting the difficulty of transfer \\
\parbox[t]{2.15cm}{Mohammadiazni\\et al. 2026} & 10 stroke patients plus healthy Ninapro DB5 pretraining data & Three-gesture classification across multiple arm postures with transfer learning & 93.6\% with transfer learning versus 46\% without transfer & Strong stroke-specific motivation for healthy-to-impaired transfer learning \\
\parbox[t]{2.15cm}{Sri-Iesaranusorn\\et al. 2021} & Ninapro DB5/DB7, mostly healthy-subject settings & 41-movement deep neural network classification & 93.87\% on DB5 and 91.69\% on DB7 & Healthy-data ceiling is high, but the task is easier than impaired-arm multilabel finger decoding \\
This project & \physiomio\ impaired-arm split from 48 stroke patients; patient-level train/val/test split & Five-finger multilabel intent decoding and four-channel hardware-targeted retraining & CNN-Large: 0.593 subset accuracy and 0.714 macro F1; distilled four-channel CNN-Micro: $0.522\pm0.011$ and $0.609\pm0.006$ & Connects impaired-arm model comparison to a concrete reduced-sensor interface \\
\bottomrule
\end{tabularx}
}
\end{table}
